\begin{theorem}[规范体积密度的高阶碰撞矩支配律]\label{thm:conclusion-gauge-density-highmoment-dominance}
\leanverified{paper\_conclusion\_gauge\_density\_highmoment\_dominance}
沿用定理 \ref{thm:conclusion-binfold-kappa-gauge-affine-law} 与命题 \ref{prop:conclusion-capacity-defect-bernstein-mellin} 的记号。对任意固定整数 \(q\ge 2\)，记
\[
S_q(m):=\sum_{x\in X_m}d_m(x)^q.
\]
则当 \(m\to\infty\) 时成立
\[
g_m
\le
\frac{1}{q-1}\log\frac{S_q(m)}{2^m}-1
+O\!\left(m\left(\frac{\varphi}{2}\right)^m\right).
\]
若进一步假设压力极限
\[
P_q:=\lim_{m\to\infty}\frac1m\log S_q(m)
\]
存在，则有斜率上界
\[
\limsup_{m\to\infty}\frac{g_m}{m}
\le
\frac{P_q-\log 2}{q-1}.
\]
因此，规范体积密度并不是独立于碰撞矩塔的额外对象；它始终被任意一个固定高阶碰撞矩所控制。
\end{theorem}
\begin{proof}
由定理 \ref{thm:conclusion-binfold-kappa-gauge-affine-law}，
\[
g_m=\kappa_m-1+O\!\left(m\left(\frac{\varphi}{2}\right)^m\right),
\qquad
\kappa_m=\sum_{x\in X_m}\mu_m(x)\log d_m(x),
\]
其中
\[
\mu_m(x):=\frac{d_m(x)}{2^m}.
\]
对随机变量 \(D_m:=d_m(X)\)（\(X\sim\mu_m\)）应用 \(\log\) 的凹性，得到
\[
\EE_{\mu_m}[\log D_m]
\le
\frac{1}{q-1}\log \EE_{\mu_m}[D_m^{\,q-1}].
\]
而
\[
\EE_{\mu_m}[D_m^{\,q-1}]
=
\sum_{x\in X_m}\frac{d_m(x)}{2^m}\,d_m(x)^{q-1}
=
\frac{S_q(m)}{2^m}.
\]
故
\[
\kappa_m\le \frac{1}{q-1}\log\frac{S_q(m)}{2^m}.
\]
与 \(g_m=\kappa_m-1+O(m(\varphi/2)^m)\) 合并，即得第一式。

若 \(P_q\) 存在，则
\[
\log\frac{S_q(m)}{2^m}
=(P_q-\log 2)m+o(m).
\]
将其代回上式并除以 \(m\)，便得
\[
\limsup_{m\to\infty}\frac{g_m}{m}
\le
\frac{P_q-\log 2}{q-1}.
\]
证毕。
\end{proof}

\begin{corollary}[规范体积的压力上包络原理]\label{cor:conclusion-gauge-density-pressure-envelope}
若对每个整数 \(q\ge 2\) 压力极限
\[
P_q=\lim_{m\to\infty}\frac1m\log S_q(m)
\]
都存在，并定义
\[
\Gamma(q):=\frac{P_q-\log 2}{q-1},
\qquad q\ge 2,
\]
则有
\[
\limsup_{m\to\infty}\frac{g_m}{m}
\le
\inf_{q\ge 2}\Gamma(q).
\]
换言之，规范体积的线性增长率不能超过碰撞矩压力谱所生成的最佳 secant envelope；若某个 \(q_\ast\) 使 \(\Gamma(q)\) 取得极小值，则该 \(q_\ast\) 给出规范体积的最紧高阶矩控制。
\end{corollary}
\begin{proof}
定理 \ref{thm:conclusion-gauge-density-highmoment-dominance} 对每个固定 \(q\ge 2\) 都给出
\[
\limsup_{m\to\infty}\frac{g_m}{m}\le \Gamma(q).
\]
对 \(q\) 取下确界即可。证毕。
\end{proof}

\begin{theorem}[Rényi--\(2\) 有效支撑的严格塌缩]\label{thm:conclusion-renyi2-effective-support-collapse}
\leanverified{paper\_conclusion\_renyi2\_effective\_support\_collapse}
记
\[
p_m(r):=\frac{d_m(r)}{2^m},
\qquad
r\in X_m,
\]
并定义 Rényi--\(2\) 有效支撑数
\[
N_{\mathrm{eff}}^{(2)}(m)
:=
\frac{1}{\sum_{r\in X_m}p_m(r)^2}
=
\frac{1}{\mathrm{Col}_m}.
\]
则有
\[
\frac{N_{\mathrm{eff}}^{(2)}(m)}{|X_m|}
=
\frac{1}{F_{m+2}\,\mathrm{Col}_m}
\longrightarrow 0.
\]
换言之，尽管稳定类型总数满足 \(|X_m|=F_{m+2}\)，在 Rényi--\(2\) 口径下真正参与统计的有效类型数只占一个趋于零的比例；可见层发生严格的 participation collapse。
\end{theorem}
\begin{proof}
由 Zeckendorf 范围上的稳定类型计数，
\[
|X_m|=F_{m+2}.
\]
另一方面，推论 \ref{cor:conclusion-collision-scale-no-single-constant-renormalization} 已给出
\[
F_{m+2}\,\mathrm{Col}_m\longrightarrow+\infty.
\]
于是
\[
\frac{N_{\mathrm{eff}}^{(2)}(m)}{|X_m|}
=
\frac{1/\mathrm{Col}_m}{F_{m+2}}
=
\frac{1}{F_{m+2}\,\mathrm{Col}_m}
\longrightarrow 0.
\]
证毕。
\end{proof}

\begin{corollary}[Rényi--\(2\) 意义下不存在渐近均匀化子列]\label{cor:conclusion-renyi2-no-asymptotic-uniformization-subsequence}
不存在任何子列 \(m_j\to\infty\) 使得分布 \(p_{m_j}\) 在 Rényi--\(2\) 口径下渐近均匀于 \(X_{m_j}\)。等价地，
\[
N_{\mathrm{eff}}^{(2)}(m_j)\sim |X_{m_j}|
\]
不可能成立。
\end{corollary}
\begin{proof}
若存在这样的子列，则有
\[
\frac{N_{\mathrm{eff}}^{(2)}(m_j)}{|X_{m_j}|}\longrightarrow 1,
\]
这与定理 \ref{thm:conclusion-renyi2-effective-support-collapse} 矛盾。证毕。
\end{proof}
